\documentclass[a4paper,12pt]{scrartcl}
\usepackage[utf8]{inputenc}
\usepackage[T1]{fontenc}
\usepackage[ngerman]{babel}
\usepackage{amsmath}
\usepackage{amssymb}
\usepackage{enumitem}
\usepackage{geometry}
\usepackage{siunitx}
\usepackage{booktabs}
\usepackage{array}
\usepackage{tikz}
\usepackage{listings}
\usepackage{xcolor}
\usepackage{scrlayer-scrpage}
\usepackage{qrcode}
\usepackage{microtype}
\usepackage[hidelinks]{hyperref}
\usetikzlibrary{arrows.meta,positioning}
\sisetup{locale=DE, per-mode=fraction, output-decimal-marker={,}}
\geometry{a4paper, left=2cm, right=2cm, top=2.1cm, bottom=2.4cm, footskip=0.9cm}

\ohead{BRAK}
\chead{Fachpraxis}
\ihead{12NP}

\RedeclareSectionCommand[beforeskip=-1.2ex, afterskip=0.5ex]{subsection}
\RedeclareSectionCommand[beforeskip=-1.0ex, afterskip=0.4ex]{subsubsection}

\lstdefinestyle{arduino}{
  language=C,
  basicstyle=\ttfamily\small,
  keywordstyle=\color{blue}\bfseries,
  commentstyle=\color{gray}\itshape,
  stringstyle=\color{red},
  numbers=none,
  frame=single,
  backgroundcolor=\color{white},
  breaklines=true,
  showstringspaces=false,
  morekeywords={HIGH,LOW,OUTPUT,INPUT,INPUT_PULLUP,pinMode,digitalWrite,digitalRead,analogRead,analogWrite,micros,millis,map,constrain,Serial,begin,println,print,delay,setup,loop,int,long,unsigned,float,void,for,while,if,else,const,A0,A1,A2},
  tabsize=2,
  xleftmargin=4pt,
  xrightmargin=4pt,
  aboveskip=3pt,
  belowskip=3pt
}

\setlength{\abovedisplayskip}{6pt}
\setlength{\belowdisplayskip}{6pt}
\setlength{\abovedisplayshortskip}{3pt}
\setlength{\belowdisplayshortskip}{3pt}
\renewcommand{\arraystretch}{1.3}
\linespread{0.97}

\newcommand{\aufg}[1]{\par\vspace{5pt}\noindent\textbf{#1}\par\nopagebreak\vspace{2pt}}

\begin{document}

\begin{center}
{\Large\textbf{Arbeitsblatt 8: Wir bauen uns ein CASSY}}\\[4pt]
\normalsize Ein eigenes Messgerät für die Entladekurve eines Kondensators \quad|\quad 16.\,Oktober\,2026
\end{center}

\vspace{2pt}
\hrule
\vspace{6pt}

\noindent
\begin{minipage}[c]{2.4cm}
\qrcode[height=2.1cm]{https://tinyurl.com/ab8-rc-sketch}
\end{minipage}%
\hspace{8pt}%
\begin{minipage}[c]{13.3cm}
\small\textbf{Der Sketch zum Herunterladen} -- diesmal wird nicht getippt:\\
\texttt{tinyurl.com/ab8-rc-sketch}\\[2pt]
Datei speichern, in der Arduino-IDE öffnen, hochladen. Lesen müssen Sie ihn trotzdem --
in Teil 2 prüfen Sie nach, ob die eingestellten Zeiten zu Ihrer Rechnung passen.
\end{minipage}

\vspace{8pt}

\noindent In Physik hängen Sie den Kondensator an das CASSY und bekommen die Kurve geschenkt.
Hier drehen wir das um: \textbf{Das Messgerät ist die Aufgabe, der Kondensator nur der Prüfling.}
Und Sie messen etwas, das mit Stoppuhr und Zeigerinstrument \emph{prinzipiell} nicht mehr geht.

\subsection*{Teil 1 -- Vorüberlegung: erst rechnen, dann aufbauen}

\noindent Gemessen wird ein Folienkondensator mit $C=\SI{100}{\nano\farad}$ über einen Widerstand
$R=\SI{47}{\kilo\ohm}$. Für Laden und Entladen gilt
\[
  \text{laden: } u(t) = U_0\left(1-e^{-t/\tau}\right), \qquad
  \text{entladen: } u(t) = U_0\,e^{-t/\tau}, \qquad \tau = R\cdot C .
\]

\aufg{Aufgabe 1 -- Wann ist der Kondensator voll, wann leer?}
Berechnen Sie $\tau$. Berechnen Sie dann, welcher Anteil von $U_0$ beim \emph{Laden} nach
$1\tau$, $3\tau$, $5\tau$ und $7\tau$ erreicht ist, und welcher Anteil beim \emph{Entladen}
jeweils noch übrig ist. \\
Legen Sie damit fest, ab wann Sie den Kondensator „voll geladen“ nennen -- und begründen Sie
Ihr Kriterium mit der Auflösung des Wandlers aus Arbeitsblatt 4: Was ist die kleinste Spannung,
die der Arduino überhaupt noch von „ganz voll“ unterscheiden kann?

\aufg{Aufgabe 2 -- Wie lange muss das Gerät messen, wie schnell muss es abtasten?}
Berechnen Sie die Halbwertszeit $t_{1/2}=\tau\cdot\ln 2$.
Die Messung soll etwa $4\tau$ lang sein und dabei $N=100$ Messpunkte aufnehmen
(das sind 99 Zeitintervalle). Berechnen Sie den nötigen Punktabstand \texttt{DT}
und runden Sie auf einen glatten Wert.
Rechnen Sie außerdem nach, ob die im Sketch eingestellte Ladezeit von \SI{500}{\milli\second}
nach Ihrem Kriterium aus Aufgabe 1 ausreicht.

\aufg{Aufgabe 3 -- Warum überhaupt ein Gerät?}
Im Physik-Praktikum messen Sie dieselbe Kurve mit $C=\SI{100}{\micro\farad}$ und demselben
Widerstand. Berechnen Sie $\tau$ und $t_{1/2}$ für diesen Fall und vergleichen Sie mit Ihren
Werten aus Aufgabe 2. Eine Reaktionszeit beim Drücken der Stoppuhr beträgt etwa
\SI{0{,}2}{\second}. Beurteilen Sie damit beide Fälle: Wo ist Handmessung möglich, wo nicht?

\subsection*{Teil 2 -- Das Gerät aufbauen und prüfen}

\noindent
\begin{minipage}[c]{9.2cm}
\begin{center}
\begin{tikzpicture}[line width=0.9pt, scale=1.0]
  \draw (0,0) rectangle (2.5,3.2);
  \node[align=center] at (1.25,1.6) {Arduino\\Uno};
  \draw (2.5,2.7) -- (4.0,2.7);
  \draw (5.2,2.7) -- (6.6,2.7) -- (6.6,1.7);
  \draw (2.5,1.7) -- (6.6,1.7);
  \draw (6.6,1.7) -- (6.6,1.15);
  \draw (6.6,0.55) -- (6.6,0.4) -- (2.5,0.4);
  \draw[fill=white] (4.0,2.45) rectangle (5.2,2.95);
  \node at (4.6,2.7) {$R$};
  \draw (6.25,1.15) -- (6.95,1.15);
  \draw (6.25,0.85) -- (6.95,0.85);
  \draw (6.6,0.85) -- (6.6,0.55);
  \node[left] at (6.15,1.0) {$C=\SI{100}{\nano\farad}$};
  \fill (6.6,1.7) circle (2.2pt);
  \node[above] at (3.1,2.7) {\small D2};
  \node[above] at (3.1,1.7) {\small A0};
  \node[above] at (3.1,0.4) {\small GND};
\end{tikzpicture}
\end{center}
\end{minipage}%
\hfill
\begin{minipage}[c]{7.4cm}
\small Der Arduino macht hier \emph{beides}: Über \texttt{D2} lädt und entlädt er den Kondensator,
über \texttt{A0} misst er dessen Spannung. Er weiß deshalb \textbf{exakt}, wann $t=0$ ist --
er hat den Moment ja selbst ausgelöst. Ein Startsignal von außen braucht er nicht.

\smallskip
Material: Arduino Uno, Steckbrett, $C=\SI{100}{\nano\farad}$ Folie,
$R=\SI{47}{\kilo\ohm}$, Multimeter.
\end{minipage}

\aufg{Aufgabe 4 -- Aufbauen, messen, Zeitbasis prüfen}
Messen Sie den Widerstand vor dem Einbau mit dem Multimeter nach und notieren Sie den
\textbf{Istwert} -- mit dem rechnen Sie später, nicht mit dem Aufdruck. Bauen Sie die Schaltung
auf, laden Sie den Sketch hoch und öffnen Sie den Seriellen Monitor mit \textbf{115200\,Baud}.\\
Vergleichen Sie den eingestellten Takt \texttt{DT} mit Ihrem Ergebnis aus Aufgabe 2.
Notieren Sie aus der Kopfzeile die \emph{Soll-} und die \emph{Ist-Messdauer} und erklären Sie
die Abweichung. Was würde eine stark anwachsende Ist-Dauer bedeuten?

\aufg{Aufgabe 5 -- Zwei Grenzen des Geräts}
\begin{enumerate}[label=(\alph*), leftmargin=2em, itemsep=2pt, topsep=2pt, partopsep=0pt, parsep=0pt]
  \item Eine einzelne Wandlung mit \texttt{analogRead()} dauert \SI{112}{\micro\second}.
        Prüfen Sie, ob der eingestellte Takt zulässig ist. Berechnen Sie dann, welche
        Kapazität Sie mit $R=\SI{47}{\kilo\ohm}$ mindestens brauchen, damit noch
        40~Messpunkte auf die Kurve passen. Wäre ein \SI{10}{\nano\farad}-Kondensator
        noch messbar?
  \item Das Gerät speichert alle Werte erst und sendet danach. Eine Messzeile ist etwa
        10~Zeichen lang, pro Zeichen werden 10~Bit übertragen. Berechnen Sie die Sendedauer
        \emph{einer} Zeile bei 115200\,Baud und vergleichen Sie sie mit dem Messtakt.
        Begründen Sie damit die Reihenfolge „erst messen, dann senden“.
\end{enumerate}

\subsection*{Teil 3 -- Auswerten}

\aufg{Aufgabe 6 -- Kurve und Schnellschätzung}
Messen Sie die Spannung am \texttt{5V}-Pin und geben Sie die Umrechnung $u(\text{adc})$ an.
Übertragen Sie die Messwerte in die Tabellenkalkulation und zeichnen Sie $u$ über $t$.
Lesen Sie die Halbwertszeit ab und vergleichen Sie mit Ihrer Rechnung aus Aufgabe 2.

\aufg{Aufgabe 7 -- Linearisierung}
Legen Sie eine Hilfsspalte $\ln u$ an und tragen Sie $\ln u$ über $t$ in ein
\textbf{normales, linear geteiltes} Achsenkreuz ein. Begründen Sie mit der Entladeformel,
warum sich eine Gerade mit der Steigung $m=-1/\tau$ ergeben muss.
Entscheiden Sie \emph{vor} dem Zeichnen der Ausgleichsgeraden, welche Messpunkte Sie
verwenden -- schauen Sie sich dazu die letzten Zeilen Ihrer Messreihe an.
Bestimmen Sie $\tau$ mit einem Steigungsdreieck.

\aufg{Aufgabe 8 -- Kapazität}
Berechnen Sie $C=\tau/R$ mit dem gemessenen Widerstandswert und geben Sie die Abweichung
vom Aufdruck \SI{100}{\nano\farad} in Prozent an.

\subsection*{Teil 4 -- Wie gut ist unser Messgerät?}

\aufg{Aufgabe 9 -- Fehlerbilanz}
\begin{enumerate}[label=(\alph*), leftmargin=2em, itemsep=2pt, topsep=2pt, partopsep=0pt, parsep=0pt]
  \item Zeigen Sie rechnerisch: Für die \emph{Steigung} der Geraden ist es gleichgültig, ob
        Sie $\ln u$ oder $\ln(\text{adc})$ auftragen. Welche überraschende Folgerung ergibt
        sich daraus für Ihr Messgerät -- im Unterschied zum Füllstandsgeber aus Arbeitsblatt 7?
  \item Die Zeitachse des Arduino ist in Schritten von \SI{4}{\micro\second} aufgelöst, der
        Taktgeber des Boards geht aber nur auf \SI{0{,}5}{\percent} genau. Rechnen Sie beide
        Beiträge in einen Fehler von $\tau$ um. Vergleichen Sie mit der Genauigkeit, mit der
        Sie den Widerstand gemessen haben.
  \item Beurteilen Sie damit: Ist die Abweichung aus Aufgabe 8 ein Messfehler Ihres Geräts
        oder eine Eigenschaft des Kondensators?
\end{enumerate}

\aufg{Aufgabe 10 (Zusatz) -- Probe mit einem anderen Widerstand}
Wiederholen Sie die Messung mit $R=\SI{470}{\kilo\ohm}$. Berechnen Sie vorher das neue $\tau$
und das passende \texttt{DT} und ändern Sie im Sketch \emph{nur diese eine Zeile}; prüfen Sie
auch, ob die Ladezeit noch reicht. Sie erhalten eine ganz andere Kurve -- aber dieselbe
Kapazität. Erklären Sie, warum das eine starke Bestätigung der Methode ist.

\end{document}
